\chapter{Exact Galerkin as Finite Compilation}
\label{chap:exact-galerkin-finite-compilation}

Galerkin computation is the place where the weak contract becomes a finite compilation. Choose a finite
test space
\[
  V_N = \operatorname{span}\{\phi_1,\ldots,\phi_N\}.
\]
Choose an approximate state \(u_N\) in the corresponding trial space. The weak equation is then asked only
against the selected tests:
\[
  D\mathcal A(u_N)[\phi_i] = 0,
  \qquad i=1,\ldots,N .
\]
That finite family is not an approximation to a proof in the compiler's eyes. It is a proof of a finite
statement. Approximation enters when the finite statement is related back to the full space by density,
stability, and error estimates. The compiler can keep those layers apart: exact row closure first,
analytic convergence second.

\section{Assembly as elaboration}
\label{se:assembly-as-elaboration}

Assembly is elaboration with indices. Each basis function contributes a column to the trial space and a row
to the test space. Each bilinear or nonlinear term is expanded into entries. For a linearized problem, the
Jacobian or stiffness matrix has the familiar form
\[
  K_{ij} = D^2\mathcal A(u_N)[\phi_j,\phi_i],
\]
and the residual vector has entries
\[
  r_i = D\mathcal A(u_N)[\phi_i].
\]
The compiler-facing content is that every \(K_{ij}\) and \(r_i\) is a typed scalar produced from declared
data. If coefficients are rational, symbolic, or otherwise exact, the resulting finite computation is
exact. If numerical quadrature is introduced, the quadrature rule is a new assumption in the ledger, not an
invisible property of Galerkin itself.

This is the same discipline that governed the rest of the compiler. A result is not more exact because it
sounds analytic, and not less exact because it is finite. Exactness depends on whether the object stated is
the object proved. The finite Galerkin system is a finite theorem. It says that a selected residual vector
vanishes, or has a stated norm, or satisfies a stated estimate. When convergence data are added, the theorem
connects the finite residual to the universal residual. The compiler should be asked for the theorem it can
actually close.

\section{Running forever and paying a finite cost}
\label{se:running-forever-paying-finite-cost}

The paired-particle apparatus introduced a useful pattern: the program may be capable of running forever,
but a measurement names the cost it is willing to spend. Galerkin computation has the same shape. One can
keep enriching the space,
\[
  V_1 \subset V_2 \subset V_3 \subset \cdots ,
\]
and the process has no natural final rung unless a tolerance, theorem, or external budget supplies one. The
compiler does not need the infinite ladder to terminate in order to certify a rung. It needs the rung
number, the basis, the assembly rules, and the residual reading at that rung. A run can be open-ended while
each receipt remains finite.

This matters for the weak Dirac apparatus because spinorial and gauge-coupled systems are not gentle
targets. The algebra is exact, but the search space can grow without mercy. The right compiler contract
therefore separates the eternal process from the finite certificate. A user may ask for one more basis
function, one more refinement, one more coupling term, or one more derivative row. Each request produces a
new finite compilation. None of those receipts pretends to be the limit unless the convergence theorem is
also carried.

\section{The compiler's receipt}
\label{se:compiler-galerkin-receipt}

The receipt for an exact Galerkin run has a small number of fields. It names the state space and test
space. It names the universal tensor that supplies the residual. It names the Sobolev norm or finite Gram
matrix used to measure that residual. It names the cost ceiling, if the run is bounded by one. It reports
the row residuals and the proof obligations discharged for those rows. If the residual vector is zero, the
receipt says so exactly. If the residual has a norm, the receipt says in which norm. If an error estimate
connects the finite receipt to the continuum problem, that estimate appears as another checked field.

That receipt is the compiler version of the weak solution story. The Dirac gate supplies the local typed
motion. The universal tensor supplies the residual. The first Frechet variation turns least activity into
a family of test obligations. Galerkin selection makes a finite family of those obligations available for
exact computation. The compiler neither weakens nor mystifies the mathematics. It records what was selected,
what was checked, what it cost, and which stronger conclusions still require their own laws.

\section{Back to the ledger}
\label{se:back-to-ledger}

The ledger does not reopen when the weak apparatus is added. It becomes more explicit about what kind of
receipt analytic computation produces. There is still no hidden choice in a named basis. There is still no
unbounded search masquerading as a proof when a cost ceiling is declared. There is still no continuum
theorem hidden inside a finite residual vector. What changes is that the compiler now has a clean language
for the most delicate experimental claim: the weak solution to the principle of least activity is the first
Frechet variation of the universal tensor, expressed in the Sobolev norm, and exact Galerkin computation is
the finite compilation of that weak statement.

At that point the three volumes speak through the same apparatus. The expository construction names the
measurement. The experimental construction performs the weak computation. The compiler construction says
what the machine can check about the computation and how the receipt should be read. That division is not a
loss of unity. It is the unity becoming auditable: one object, read as physics, as experiment, and as
elaboration.
